\begin{abstract}
Recent work (Falzon and Tang USENIX 2025) has shown that a protocol participant who behaves honestly but strategically chooses its inputs can break input privacy in the Private Join and Compute functionality.
In this work, we expand our understanding of attacks in this setting by investigating a broader class of functionalities, namely: Private Set Union (PSU), PSU-Cardinality (PSU-CA), and Meta's multi-key private matching (MKPM) functionality.

We first present an attack on PSU that reconstructs the intersection using only two protocol invocations. We then show that any attack on PSI-Cardinality, such as that of Guo et al.\ (USENIX~2023), lifts to an attack on PSU-CA that recovers the intersection using only three additional queries. For the MKPM protocol, we distinguish its intended matching functionality from the additional protocol-specific leakage. We give an attack against the intended functionality, and then show that exploiting the additional leakage enables even stronger attacks, including partially reconstructing the other party’s records from a single protocol invocation.

We conclude by discussing possible mitigations for deploying such systems. Our analysis demonstrates limitations of existing secure multi-party computation security definitions and highlights the real-world privacy risks associated with deploying these functionalities in practice.
\end{abstract}
